\documentclass[11pt,a4paper]{article}
\usepackage[T1]{fontenc}
\usepackage[utf8]{inputenc}
\usepackage{lmodern}
\usepackage{amsmath,amssymb}
\usepackage{bm}
\usepackage{booktabs}
\usepackage{tabularx}
\usepackage{array}
\usepackage[margin=2.5cm]{geometry}
\usepackage[colorlinks=true,linkcolor=blue,citecolor=blue,urlcolor=blue]{hyperref}

\newcommand{\vv}{\bm{v}}
\newcommand{\Jm}{\mathcal{J}}
\newcommand{\dd}{\mathrm{d}}
\newcommand{\code}[1]{\texttt{#1}}
\newcolumntype{L}{>{\raggedright\arraybackslash}X}

\title{The continuous Galerkin spectral element dynamical core\\
of CGDycore.jl (\code{src/CGSEM})\\[0.3em]
\large Numerical scheme and comparison with HOMME and ClimaCore}
\author{Summary of the implementation}
\date{\today}

\begin{document}
\maketitle

\begin{abstract}
This note describes the numerical scheme implemented in the module
\code{CGSEM} of CGDycore.jl as it is coded in the kernels of
\code{src/CGSEM}. The scheme combines a continuous Galerkin spectral element
method (SEM) on Gauss--Lobatto--Legendre (GLL) points in the horizontal with
a low-order, staggered (Charney--Phillips type) discretization in a
terrain-following height coordinate in the vertical. The prognostic variables
are density, the velocity components in vector-invariant form, the vertical
velocity on layer interfaces, the density-weighted potential temperature and
tracer densities. Time integration is done by Rosenbrock-type W-methods with an
approximate Jacobian that contains only the vertical acoustic and buoyancy
coupling (HEVI). We describe the discrete operators term by term and compare
the design with the spectral element cores HOMME (E3SM) and ClimaCore
(CliMA).
\end{abstract}

\tableofcontents

%==============================================================================
\section{Overview}
%==============================================================================

The main routine for the right-hand side is \code{Fcn!} in
\code{CGSEM/Fcn.jl}. One evaluation consists of
\begin{enumerate}
\item a diagnostic step \code{FcnPrepare!} (pressure, temperature, eddy
  coefficients),
\item a first pass of horizontal hyperviscosity (Laplacians of momentum,
  potential temperature and tracers), followed by direct stiffness summation
  (DSS) and parallel exchange,
\item the second pass, in which all tendencies are computed: biharmonic
  hyperviscosity, rotational and Coriolis terms, gradient of kinetic energy,
  pressure gradient and gravity, transport of $\rho$, $\rho\theta$ and the
  tracers,
\item DSS/parallel exchange of the tendency and the column physics (vertical
  diffusion, surface fluxes, forcing, microphysics, sedimentation, damping).
\end{enumerate}
All element-wise operators are evaluated in KernelAbstractions kernels, one
thread per horizontal node and layer. The contributions of the elements are
accumulated into the global degrees of freedom by atomic additions and divided
by the assembled lumped mass. In this way DSS is part of every operator
kernel. The elements are ordered so that the elements at the boundary of a
process come first: they are computed first, the exchange is started, the
interior elements are computed and finally the exchange is completed (overlap
of communication and computation).

\begin{table}[h]
\centering
\begin{tabularx}{\textwidth}{lL}
\toprule
Horizontal & CG spectral elements on quadrilaterals (cubed sphere or
  Cartesian), GLL collocation, polynomial order $p$ (typically $p=3$ or $4$),
  lumped mass, DSS\\
Vertical & terrain-following height coordinate, $N_z$ layers, lowest-order
  staggering: $\rho,u,v,\rho\theta,\rho q$ in layers, $w$ on interfaces\\
Momentum & vector-invariant form, prognostic velocities $u,v,w$\\
Thermodynamics & $\rho\theta$ (dry); variants with internal or total energy\\
Horizontal transport & centred SEM operators; split form for $\rho\theta$\\
Vertical transport & limited third-order upwind reconstruction
  (Koren-type limiter, non-uniform grid)\\
Dissipation & biharmonic hyperviscosity, separate coefficients for
  vorticity, divergence and scalars\\
Tracer limiter & optional optimization-based limiter (local bounds,
  mass conservation)\\
Time integration & Rosenbrock W-methods (e.g.\ \code{SSP-Knoth}) with a
  frozen vertical acoustic Jacobian, Schur complement for $w$, tridiagonal
  solves per column\\
\bottomrule
\end{tabularx}
\caption{Main components of the CGSEM scheme.}
\end{table}

%==============================================================================
\section{Continuous equations}
%==============================================================================

The dry compressible Euler equations are solved in the vector-invariant form
\begin{align}
  \partial_t \rho + \nabla\cdot(\rho\vv) &= 0, \label{eq:cont}\\
  \partial_t \vv + (\bm{\omega} + 2\bm{\Omega})\times\vv
    + \nabla K + \frac{1}{\rho}\nabla p + \nabla\Phi &= \bm{D}_{\vv},
    \label{eq:mom}\\
  \partial_t (\rho\theta) + \nabla\cdot(\rho\theta\vv) &= D_{\theta},
    \label{eq:theta}\\
  \partial_t (\rho q_l) + \nabla\cdot(\rho q_l\vv) &= D_{q_l} + S_{q_l},
    \qquad l = 1,\dots,N_{Tr}, \label{eq:tracer}
\end{align}
with $\bm{\omega} = \nabla\times\vv$, the kinetic energy
$K = \tfrac12|\vv|^2$, the geopotential $\Phi$ (gravity
$\bm g = -\nabla\Phi$, shallow or deep), the dissipation terms $\bm D$ and
the physical sources $S$. For the dry state the pressure is
\begin{equation}
  p = p_0\left(\frac{R_d\,\rho\theta}{p_0}\right)^{1/(1-\kappa)},
  \qquad \kappa = R_d/c_{pd}.
\end{equation}
The velocity $\vv = (u,v,w)$ is stored in local Cartesian components:
on the sphere $u$ and $v$ are the eastward and northward components and $w$
is the radial component. Other options select internal or total energy instead of
$\rho\theta$ (\code{DryInternalEnergy}, \code{DryTotalEnergy}, \dots), moisture
variables, a turbulent kinetic energy and EDMF variables; they reuse the same
operators and are not discussed here.

%==============================================================================
\section{Grid, coordinate transformation and metric terms}
%==============================================================================

\subsection{Horizontal elements and vertical layers}

The horizontal mesh consists of quadrilateral elements $\Omega_e$, for
example the elements of an equiangular or equidistant cubed sphere with
$n_{\text{Panel}}\times n_{\text{Panel}}$ elements per panel. Every column of
elements is divided into $N_z$ layers between the interfaces
$z_{1/2} < z_{3/2} < \dots < z_{N_z+1/2} = H$ (optionally stretched). A
grid cell $\Omega_{e,k}$ is mapped from the reference cell
$(\xi^1,\xi^2,\zeta)\in[-1,1]^3$ by
\begin{equation}
  \bm{x}(\xi^1,\xi^2,\zeta) = \frac{\bm{X}(\xi^1,\xi^2)}{|\bm{X}(\xi^1,\xi^2)|}
    \bigl(a + h(z(\zeta), z_s(\xi^1,\xi^2))\bigr),
  \quad
  z(\zeta) = \tfrac12\bigl((1-\zeta) z_{k-1/2} + (1+\zeta) z_{k+1/2}\bigr),
\end{equation}
where $\bm X$ is the bilinear interpolation of the four element corners on the
cube, $a$ is the radius of the planet and $h(z,z_s)$ is the terrain-following
transformation (\code{AdaptGrid}, e.g.\ Gal-Chen or SLEVE) of the
orography $z_s$. In the Cartesian case $\bm x$ is the analogous mapping
without the projection onto the sphere.

\subsection{Nodes and derivative matrices}

In the horizontal the GLL points $\xi_i$ with the weights $\omega_i$,
$i=1,\dots,N$, $N=p+1$, are used as interpolation and quadrature nodes. The
strong and weak derivative matrices are
\begin{equation}
  D_{ij} = \ell_j'(\xi_i), \qquad
  D^W = -W^{-1} D^T W, \quad W = \operatorname{diag}(\omega_1,\dots,\omega_N),
\end{equation}
(\code{DS} and \code{DW} in the code). In the vertical each layer has the
two GLL points $\zeta = -1$ (bottom, index $K=1$) and $\zeta=+1$ (top,
$K=2$); the corresponding two-point quadrature has the weights $1,1$.
The metric terms are stored at the $N\times N\times 2$ points of every cell.

\subsection{Metric terms}

With the Jacobian matrix $\partial\bm x/\partial\bm\xi$ in the local basis
(east, north, radial), the code stores the pointwise determinant and the
adjugate,
\begin{equation}
  J = \det\frac{\partial\bm x}{\partial\bm\xi}, \qquad
  \Jm^{a}_{\;b} = J\,\frac{\partial\xi^a}{\partial x^b},
  \qquad a,b = 1,2,3,
\end{equation}
(\code{J} and \code{dXdxI[a,b,K,ID,k,e]}, $\xi^3=\zeta$). The rows
$\Jm^a = J\nabla\xi^a$ are the contravariant basis vectors scaled by $J$. The
derivative of the height with respect to $\zeta$ is computed by the two-point
difference $\partial_\zeta h \approx \tfrac12(h_{K=2}-h_{K=1})$ and the
horizontal derivatives of $h$ by the matrix $D$. The divergence, gradient and
curl are written in the metric-conservative forms
\begin{equation}
  J\nabla\cdot\bm f = \sum_{a}\partial_{\xi^a}\bigl(\Jm^a\cdot\bm f\bigr),\quad
  J\nabla\phi = \sum_a \Jm^a\,\partial_{\xi^a}\phi,\quad
  J(\nabla\times\vv)_i = \sum_a \varepsilon_{ibc}\,
    \partial_{\xi^a}\bigl(\Jm^a_{\;b}\,v_c\bigr),
  \label{eq:metricops}
\end{equation}
which are equivalent to the usual forms because of the metric identities
$\sum_a\partial_{\xi^a}\Jm^a = 0$.

\subsection{Lumped mass matrix and direct stiffness summation}

Let $g$ be a global horizontal node and $k$ a layer. The assembled half-cell
masses are
\begin{equation}
  M^{(K)}_{g,k} = \sum_{(e,i,j)\mapsto g} J^{(K)}_{e,ij,k}, \qquad K=1,2,
  \qquad M_{g,k} = M^{(1)}_{g,k} + M^{(2)}_{g,k}
\end{equation}
(\code{MassCGKernel!}, array \code{M[k,g,K]}). The sum runs over all copies of
the node $g$ in the elements that contain it. The exact lumped GLL mass would
be $\omega_i\omega_j M_{g,k}$; on a conforming quadrilateral mesh the product
$\omega_i\omega_j$ is the same for all copies of a node (the GLL weights are
symmetric), so the weights cancel in every quotient
``assembled contribution / assembled mass'' and are omitted in both.

A generic term is evaluated in strong form element by element and assembled
as a $J$-weighted average,
\begin{equation}
  F_{g,k} = \frac{1}{M_{g,k}} \sum_{(e,i,j)\mapsto g} R_{e,ij,k},
  \label{eq:dss}
\end{equation}
where $R_{e,ij,k}$ is $J$ times the pointwise tendency, summed over the two
vertical nodes of the layer. This is the SEM collocation with lumped mass
together with DSS.
The vertical velocity lives on the interfaces $k+\tfrac12$ and gets the
assembled interface masses $M^{(2)}_{g,k}+M^{(1)}_{g,k+1}$, weighted with the
density (Section~\ref{sec:w}).

%==============================================================================
\section{Vertical staggering and boundary conditions}
%==============================================================================

Per horizontal node and layer $k=1,\dots,N_z$ the unknowns
$\rho_k, u_k, v_k, (\rho\theta)_k, (\rho q_l)_k$ are constant in the layer.
The vertical velocity $w_{k+1/2}$, $k=1,\dots,N_z-1$, is stored on the interior
interfaces (array position $k$ belongs to the top of layer $k$). Inside a
layer the velocity at the two vertical nodes is
\begin{equation}
  \vv^{(1)}_k = (u_k, v_k, w_{k-1/2}), \qquad \vv^{(2)}_k = (u_k, v_k, w_{k+1/2}),
\end{equation}
so $u,v$ are piecewise constant and $w$ is piecewise linear in $\zeta$
(Charney--Phillips staggering of $w$ against $\rho,\rho\theta$).
At the top $w_{N_z+1/2}=0$. At the lower boundary the kinematic condition
$\Jm^3\cdot\vv = 0$ (no flow through the orography) gives
\begin{equation}
  w_{1/2} = -\frac{\Jm^3_{\;1}u_1 + \Jm^3_{\;2}v_1}{\Jm^3_{\;3}}.
\end{equation}
The vertical fluxes of mass, $\rho\theta$ and tracers vanish at the bottom and
at the top.

%==============================================================================
\section{Spatial discretization of the individual terms}
%==============================================================================

In the following, sums over $m$ are horizontal SEM derivatives inside an
element, e.g.\ $(D_1 f)_{ij} = \sum_m D_{im}f_{mj}$ and
$(D_2 f)_{ij} = \sum_m D_{jm}f_{im}$. Every element contribution is assembled
with \eqref{eq:dss}.

\subsection{Continuity equation}

The horizontal mass fluxes in contravariant form, summed over both vertical
nodes of the layer, are
\begin{equation}
  U^a_k = \rho_k \sum_{K=1}^2 \bigl(\Jm^{a,(K)}_{\;1} u_k + \Jm^{a,(K)}_{\;2} v_k\bigr),
  \qquad a = 1,2,
\end{equation}
(\code{Contra12}). The vertical mass flux through the interface $k+\tfrac12$
uses the metric from both sides,
\begin{equation}
  W_{k+1/2} = \rho_k\,\Jm^{3,(2)}_k\cdot(u_k,v_k,w_{k+1/2})
    + \rho_{k+1}\,\Jm^{3,(1)}_{k+1}\cdot(u_{k+1},v_{k+1},w_{k+1/2})
\end{equation}
(\code{Contra3}). The element contribution to the density tendency is
\begin{equation}
  R^\rho_k = -\bigl(D_1 U^1_k + D_2 U^2_k\bigr)
     - \tfrac12\bigl(W_{k+1/2} - W_{k-1/2}\bigr).
\end{equation}

\subsection{Potential temperature: horizontal split form}

With $\theta_k = (\rho\theta)_k/\rho_k$ the horizontal part of the flux
divergence is discretized in the split form
\begin{equation}
  R^{\rho\theta,h}_k = -\tfrac12\Bigl(D_1(U^1\theta) + D_2(U^2\theta)
    + U^1 D_1\theta + U^2 D_2\theta + \theta\,(D_1 U^1 + D_2 U^2)\Bigr)_k
\end{equation}
(\code{DivRhoTrUpwind3New2Kernel!}). Together with the discrete continuity
equation this split form keeps a constant $\theta$ constant (the
split form reduces to $\theta$ times the mass flux divergence), and it is a
skew-symmetric (kinetic-energy-preserving type) discretization of the
advection of $\theta$. Tracers $\rho q_l$ use the plain flux form
$-(D_1(U^1 q_l) + D_2(U^2 q_l))$.

\subsection{Vertical transport of \texorpdfstring{$\rho\theta$}{rho theta} and tracers}
\label{sec:upwind}

The vertical fluxes of $\rho\theta$ and $\rho q_l$ use the same mass flux
$\tfrac12 W_{k+1/2}$ as the continuity equation (consistency with mass) and
an upwind value of the specific quantity $c = \theta$ or $q_l$,
\begin{equation}
  \mathcal{F}_{k+1/2} = \tfrac14\Bigl((|W| + W)\,c^{L}_{k+1/2}
    + (W - |W|)\,c^{R}_{k+1/2}\Bigr), \qquad W = W_{k+1/2}.
\end{equation}
The interface values $c^{L}$, $c^{R}$ are reconstructed from the four layers
$k-1,k,k+1,k+2$ with a limited third-order upwind-biased reconstruction on
the non-uniform grid (\code{RecU4}). With the layer volumes
$J_k = \sum_K J^{(K)}_k$ and
\begin{equation}
  r = \frac{c_k - c_{k-1}}{c_{k+1} - c_k}, \qquad
  \kappa_R = \frac{J_k}{J_k+J_{k+1}}\,\frac{J_{k-1}+J_k}{J_{k-1}+J_k+J_{k+1}},\quad
  \kappa_L = \frac{J_k}{J_{k-1}+J_k}\,\frac{J_{k+1}}{J_{k-1}+J_k+J_{k+1}},
\end{equation}
the left value is
\begin{equation}
  c^{L}_{k+1/2} = c_k + \max\bigl(0, \min(r,\, \kappa_L r + \kappa_R,\, 1)\bigr)\,(c_{k+1}-c_k),
\end{equation}
and $c^R$ is defined mirror-symmetrically. On a uniform grid
$\kappa_L = 1/6$, $\kappa_R = 1/3$, and without the limiter
$c^L = -\tfrac16 c_{k-1} + \tfrac56 c_k + \tfrac13 c_{k+1}$, the
third-order upwind-biased value. The limiter is of Koren type \cite{Koren1993} (TVD, bounded by
$0$ and $1$ and by the first-order slope $r$). Near the boundaries the missing
layers are replaced by the boundary layer.

\subsection{Momentum: rotational term}

The vorticity is evaluated with the conservative curl \eqref{eq:metricops}.
The horizontal derivatives are taken with $D$ separately at both vertical
nodes $K=1,2$ of the layer. The vertical derivative of a flux $f$ (for example
$\Jm^3_{\;3}v - \Jm^3_{\;2}w$) has two parts:
\begin{itemize}
\item inside the layer the two-point difference $\tfrac12(f^{(2)}_k - f^{(1)}_k)$,
\item at each interface the jump $\tfrac12(f^{(1)}_{k+1} - f^{(2)}_k)$; each of
  the two adjacent half cells gets half of it (a central interface
  treatment, comparable to a DG lifting with the central flux).
\end{itemize}
The cross product with $\vv$ is formed pointwise at the vertical nodes with the
interface values of $w$. The horizontal momentum tendency is
\begin{equation}
  \partial_t (u,v)_k \leftarrow
  -\frac{1}{M_{g,k}}\sum_{\text{copies}}\sum_{K=1}^2
  \Bigl[(J\bm\omega^{(K)})\times\vv^{(K)}\Bigr]_{1,2}
\end{equation}
(\code{MomentumVectorInvariantCoriolisKernel!}). The Coriolis force is added
in the same kernel. With the shallow-atmosphere approximation only the
vertical component $f=2\Omega\sin\varphi$ is used, multiplied by
$J^{(1)}+J^{(2)}$. Variants: $f$-plane, deep atmosphere.

\subsection{Momentum: kinetic energy, pressure gradient and gravity}

The kinetic energy at the two vertical nodes of a layer is
\begin{equation}
  K^{(1)}_k = \tfrac12\bigl(u_k^2 + v_k^2 + w_{k-1/2}^2\bigr),\qquad
  K^{(2)}_k = \tfrac12\bigl(u_k^2 + v_k^2 + w_{k+1/2}^2\bigr).
\end{equation}
Its gradient $J\nabla K = \sum_a\Jm^a\partial_{\xi^a}K$ is formed with $D$ in
the horizontal and with the same two-part vertical difference as above
(\code{GradFullKernel!}).

The horizontal pressure gradient uses the horizontal derivatives of the
layer pressure $p_k$:
\begin{equation}
  (J\nabla p)_b \approx \sum_{K=1}^2\sum_{a=1}^2 \Jm^{a,(K)}_{\;b}\,\partial_{\xi^a}p_k
  + \Bigl(\sum_{K}\Jm^{3,(K)}_{\;b}\Bigr)\,\partial_\zeta p,
  \qquad
  \partial_\zeta p \approx -g\,\rho_k\,
  \frac{J^{(1)}_k + J^{(2)}_k}{\Jm^{3,(1)}_{3,k} + \Jm^{3,(2)}_{3,k}},
  \qquad b = 1,2.
\end{equation}
Because $\Jm^3_{\;3} = J\,\partial\zeta/\partial z$, the quotient is
$\partial z/\partial\zeta$, so the vertical pressure derivative in the
terrain-following correction term is replaced by its hydrostatic value
$\partial_z p = -\rho g$. This is the treatment of the horizontal pressure
gradient over terrain used in MPAS \cite{KlempSkamarock2021}. It needs no
vertical average of $\partial_\zeta p$ from the interfaces to the layer
centres, and it is exact for hydrostatic states. The horizontal tendency is
$-(J\nabla p)_{1,2}/\rho_k$, assembled with $M_{g,k}$. Without orography
$\Jm^3_{\;1} = \Jm^3_{\;2} = 0$ and the correction term vanishes.

\subsection{Vertical momentum}
\label{sec:w}

All contributions to the vertical velocity are formed in the two half cells
that touch the interface $k+\tfrac12$ (top half of layer $k$, bottom half of
layer $k+1$) and assembled with density-weighted masses:
\begin{equation}
  \partial_t w_{k+1/2} =
  \frac{\sum_{\text{copies}}\bigl(\rho_k R^{w,(2)}_k + \rho_{k+1} R^{w,(1)}_{k+1}
    + R^{p}_{k+1/2}\bigr)}
  {\sum_{\text{copies}}\bigl(\rho_k J^{(2)}_k + \rho_{k+1} J^{(1)}_{k+1}\bigr)}.
\end{equation}
Here $R^{w,(K)}$ contains the rotational term, the vertical gradient of $K$
and the Coriolis term (deep atmosphere only). The vertical pressure gradient and
gravity are
\begin{equation}
  R^p_{k+1/2} = -\tfrac12\bigl(p_{k+1} - p_k\bigr)
     \bigl(\Jm^{3,(2)}_{3,k} + \Jm^{3,(1)}_{3,k+1}\bigr)
     - g\bigl(\rho_k J^{(2)}_k + \rho_{k+1} J^{(1)}_{k+1}\bigr),
\end{equation}
i.e.\ $\bar\rho\,\partial_t w = -\partial_z p - \bar\rho g$ with the
volume-weighted interface density $\bar\rho$. For a hydrostatic state
$(p_{k+1}-p_k)/\Delta z_{k+1/2} = -g\bar\rho_{k+1/2}$ is the discrete
balance that the right-hand side preserves exactly.

\subsection{Discrete kinetic energy and vertical averaging}
\label{sec:kin}

The definition of the discrete kinetic energy and the averaging to the
interfaces are chosen such that the vertical discretization conserves the
density-weighted kinetic energy of a column. We show this for the
two-dimensional system in the coordinates $(\xi,\zeta)$ written in covariant
($u_\xi,u_\zeta$) and contravariant ($u^\xi,u^\zeta$) components,
\begin{align}
  J\partial_t\rho &= -\partial_\xi(J\rho u^\xi) - \partial_\zeta(J\rho u^\zeta),\\
  J\rho\,\partial_t u_\xi &= -J\rho\,\partial_\xi K
     - J\rho u^\zeta\,\partial_\xi u_\zeta + J\rho u^\zeta\,\partial_\zeta u_\xi,\\
  J\rho\,\partial_t u_\zeta &= -J\rho\,\partial_\zeta K
     + J\rho u^\xi\,\partial_\xi u_\zeta - J\rho u^\xi\,\partial_\zeta u_\xi,
\end{align}
with $K = \tfrac12(u_\xi u^\xi + u_\zeta u^\zeta)$; pressure gradient and
gravity are omitted. The covariant $u_k$ (and $u^\xi$) live in the layers,
$w_{k+1/2}$ (and $u^\zeta$) on the interfaces.

\paragraph{Half-cell contravariant velocities.}
With the symmetric metric tensor $G_k$ of layer $k$, which maps covariant to
contravariant components, the velocities of the lower (L) and upper (R) half
cell are
\begin{equation}
  \begin{pmatrix} U^L_k\\ W^L_k\end{pmatrix}
    = G_k\begin{pmatrix} u_k\\ w_{k-1/2}\end{pmatrix},\qquad
  \begin{pmatrix} U^R_k\\ W^R_k\end{pmatrix}
    = G_k\begin{pmatrix} u_k\\ w_{k+1/2}\end{pmatrix}.
\end{equation}
The boundary values $w_{1/2}$ and $w_{N+1/2}$ are chosen such that
$W^L_1 = W^R_N = 0$. The layer and interface values are
\begin{equation}
  U_k = \tfrac12\bigl(U^L_k + U^R_k\bigr),\qquad
  W_{k+1/2} = \frac{J_k\rho_k W^R_k + J_{k+1}\rho_{k+1}W^L_{k+1}}
                   {J_k\rho_k + J_{k+1}\rho_{k+1}},
  \label{eq:Wavg}
\end{equation}
i.e.\ the contravariant vertical velocity at an interface is the
mass-weighted mean of the two adjacent half cells.

\paragraph{Kinetic energy.}
The kinetic energy of a layer is the mean of the two half-cell energies,
\begin{equation}
\begin{split}
  K_k &= \frac14\begin{pmatrix} u_k\\ w_{k-1/2}\end{pmatrix}^{\!T}
          G_k\begin{pmatrix} u_k\\ w_{k-1/2}\end{pmatrix}
      + \frac14\begin{pmatrix} u_k\\ w_{k+1/2}\end{pmatrix}^{\!T}
          G_k\begin{pmatrix} u_k\\ w_{k+1/2}\end{pmatrix}\\
      &= \frac14\bigl(u_k(U^L_k+U^R_k) + w_{k-1/2}W^L_k + w_{k+1/2}W^R_k\bigr).
\end{split}
  \label{eq:Kdisc}
\end{equation}
Because $G_k$ is symmetric,
$\dot K_k = \tfrac12\bigl(\dot u_k(U^L_k+U^R_k) + \dot w_{k-1/2}W^L_k
  + \dot w_{k+1/2}W^R_k\bigr)$.

\paragraph{Semi-discrete equations.}
Let $D_x$ be the horizontal derivative, $D_{k+1/2}$ a discretization of
$\partial_\zeta u_\xi$ at the interface and
$\Omega_{k+1/2} = D_{k+1/2} - D_x w_{k+1/2}$ the discrete vorticity there.
Then
\begin{align}
  J_k\dot\rho_k &= -D_x(J_k\rho_kU_k)
    - \tfrac12(J_k\rho_k + J_{k+1}\rho_{k+1})W_{k+1/2}\notag\\
    &\quad + \tfrac12(J_{k-1}\rho_{k-1} + J_k\rho_k)W_{k-1/2},
    \label{eq:rho2d}\\
  J_k\rho_k\dot u_k &= -J_k\rho_k D_xK_k
    + \tfrac12 J_k\rho_k\bigl(W_{k+1/2}\Omega_{k+1/2}
    + W_{k-1/2}\Omega_{k-1/2}\bigr),
    \label{eq:u2d}\\
  \tfrac12(J_k\rho_k + J_{k+1}\rho_{k+1})\,\dot w_{k+1/2} &=
    -\tfrac12(J_k\rho_k + J_{k+1}\rho_{k+1})(K_{k+1} - K_k)\notag\\
    &\quad - \tfrac12(J_k\rho_kU_k + J_{k+1}\rho_{k+1}U_{k+1})\,\Omega_{k+1/2}.
    \label{eq:w2d}
\end{align}
By \eqref{eq:Wavg}, the vertical mass flux in \eqref{eq:rho2d} is
$\tfrac12(J_k\rho_kW^R_k + J_{k+1}\rho_{k+1}W^L_{k+1})$.

\paragraph{Column energy.}
Summing $J_k\rho_kK_k$ over the column and using \eqref{eq:Wavg} and the
boundary conditions gives
\begin{equation}
  \mathcal K = \sum_{k=1}^N J_k\rho_kK_k
   = \sum_{k=1}^N J_k\rho_k u_kU_k
   + \sum_{k=1}^{N-1}\tfrac12(J_k\rho_k + J_{k+1}\rho_{k+1})\,w_{k+1/2}W_{k+1/2},
\end{equation}
and therefore
\begin{equation}
\begin{split}
  \dot{\mathcal K} = {}&\sum_{k=1}^N J_k\dot\rho_kK_k
   + \sum_{k=1}^N J_k\rho_k\dot u_kU_k\\
   &+ \sum_{k=1}^{N-1}\tfrac12(J_k\rho_k + J_{k+1}\rho_{k+1})\,\dot w_{k+1/2}W_{k+1/2}.
\end{split}
\end{equation}
Inserting \eqref{eq:rho2d}--\eqref{eq:w2d}, the vertical terms cancel pairwise:
\begin{itemize}
\item The vorticity terms cancel between the $u$ and the $w$ equation.
  This requires the same interface velocity $W_{k+1/2}$ in \eqref{eq:u2d} as
  in the energy, and the mass-weighted layer velocity in \eqref{eq:w2d}.
\item The vertical mass flux in the continuity equation times $K_k$ cancels,
  after summation by parts, against the vertical kinetic energy gradient
  in \eqref{eq:w2d}. This requires the same weight
  $\tfrac12(J_k\rho_k + J_{k+1}\rho_{k+1})$ in both places and the layer
  mean \eqref{eq:Kdisc} for $K$.
\end{itemize}
What remains are boundary terms proportional to $W_{1/2}$ and $W_{N+1/2}$,
which vanish. The horizontal terms are handled by the SEM summation-by-parts
property of $D_x$ with the lumped GLL mass.

\paragraph{Relation to the implementation.}
The code uses Cartesian velocity components, so $G_k$ becomes the identity in
these components, and the metric enters through $\Jm^3$:
\begin{itemize}
\item The half-cell energies $K^{(1)}_k$, $K^{(2)}_k$ of
  \code{GradFullKernel!} are the two terms in \eqref{eq:Kdisc}, i.e.\
  $K_k = \tfrac12(K^{(1)}_k + K^{(2)}_k)$. $K_k$ itself is never formed.
  The horizontal gradient is applied to each half-cell energy with the
  metric of its node and weighted with the half-cell masses.
\item The vertical mass flux of \code{Contra3} is exactly
  $\tfrac12(J_k\rho_kW^R_k + J_{k+1}\rho_{k+1}W^L_{k+1})$ with
  $J_kW^R_k = \Jm^{3,(2)}_k\cdot\vv^{(2)}_k$, i.e.\ the averaging
  \eqref{eq:Wavg}. The same flux is used for $\rho\theta$ and the tracers.
\item The rotational term is formed pointwise in each half cell with the
  Cartesian half-cell velocity. Since $(\bm\omega\times\vv)\cdot\vv = 0$ and the
  $w$ contributions are weighted with $\rho_k$ and $\rho_{k+1}$ in the
  denominator of Section~\ref{sec:w}, it is energy neutral half cell by half
  cell. This takes the role of the averaged $W_{k+1/2}$ in \eqref{eq:u2d}.
\item The vertical kinetic energy gradient in the $w$ equation is assembled
  half cell by half cell. In flat coordinates this gives
  \begin{equation}
    \tfrac12\rho_k\bigl(K^{(1)}_{k+1} - K^{(1)}_k\bigr)
    + \tfrac12\rho_{k+1}\bigl(K^{(2)}_{k+1} - K^{(2)}_k\bigr)
  \end{equation}
  instead of $\tfrac12(\rho_k + \rho_{k+1})(K_{k+1} - K_k)$ in
  \eqref{eq:w2d}. The difference is
  $\tfrac14(\rho_k - \rho_{k+1})
   \bigl[(K^{(1)}_{k+1} - K^{(1)}_k) - (K^{(2)}_{k+1} - K^{(2)}_k)\bigr]$.
  It vanishes for $\rho_k = \rho_{k+1}$ but not in general, so the column
  energy identity above holds for the code only up to this term.
\end{itemize}

\subsection{Horizontal hyperviscosity}

The dissipation $\bm D_{\vv}$, $D_\theta$, $D_{q_l}$ is a biharmonic operator
in the horizontal, built in two passes with a DSS in between
(\code{HyperViscKernel!}, \code{HyperViscKoeffKernel!}).
\begin{enumerate}
\item First pass: per layer the vertical vorticity
  $\zeta_v = \bm k\cdot\nabla\times\vv$ and the divergence
  $\delta = \nabla_h\cdot\vv$ are computed in strong form and divided by
  $J^{(1)}+J^{(2)}$. Then the weak forms of the rotational and the divergent
  part of the vector Laplacian,
  \begin{equation}
    \bm L_{\text{rot}} = -\bm k\times\nabla^W \zeta_v,\qquad
    \bm L_{\text{div}} = \nabla^W\delta,\qquad
    L_\theta = \nabla^W\cdot\nabla\theta,
  \end{equation}
  are formed with $D^W$ (so that $\bm L_{\text{rot}}+\bm L_{\text{div}}
  = \nabla(\nabla\cdot\vv) - \nabla\times(\nabla\times\vv)$ is the vector
  Laplacian). They are stored separately and assembled (DSS + exchange).
\item Second pass: the same operators are applied to the stored fields:
  \begin{equation}
    \bm D_{\vv} = -\nu_{\text{curl}}\,\Delta_{\text{rot}}\bm L_{\text{rot}}
       - \nu_{\text{grad}}\,\Delta_{\text{div}}\bm L_{\text{div}},\qquad
    D_\theta = -\nu_{\theta}\,\nabla^W\cdot\bigl(\rho\,\nabla L_\theta\bigr).
  \end{equation}
\end{enumerate}
The coefficients are \code{HyperDCurl}, \code{HyperDGrad} and
\code{HyperDDiv}. Optionally there is a hyperviscosity for $w$
(\code{HyperDDivW}). Tracers are treated like $\theta$. The default scaling
with the resolution (\code{HyperViscosity}) is
\begin{equation}
  \nu = 5\cdot 10^{14}\,\mathrm{m^4\,s^{-1}}\;
  \Bigl(\frac{32}{n_{\text{Panel}}}\Bigr)^{\log_2 10}
  \approx 5\cdot 10^{14}\,\Bigl(\frac{32}{n_{\text{Panel}}}\Bigr)^{3.32},
\end{equation}
comparable to the $\Delta x^{3.2}$ scaling that is used in HOMME.

\subsection{Optional horizontal limiter for tracers}

With \code{HorLimit=true} the horizontal tracer transport uses an
optimization-based limiter in the spirit of Guba et al.\ (2014)
(\code{LimitKernel!}, \code{DivRhoTrUpwind3LimKernel!}):
\begin{enumerate}
\item the minimum and maximum of $q$ in every element are computed, and the
  bounds $q_{\min}, q_{\max}$ are taken over the element and its neighbours
  (a stencil of 13 elements);
\item a provisional explicit step $\tilde\rho$, $\widetilde{\rho q}$ with the
  current stage time step is formed;
\item a constant shift $\lambda$ is determined per element and layer so that
  $q^\star = \operatorname{median}(q_{\min}, \widetilde{\rho q}/\tilde\rho + \lambda, q_{\max})$
  conserves the tracer mass in the element,
  $\sum \omega_i\omega_j J\,(q^\star\tilde\rho - \widetilde{\rho q}) = 0$.
  The scalar equation is solved by a secant iteration (at most 8 steps).
\end{enumerate}
The limited tendency is then derived from $q^\star$. The vertical part
uses the reconstruction of Section~\ref{sec:upwind}.

%==============================================================================
\section{Time integration}
%==============================================================================

\subsection{Rosenbrock W-methods}

The semi-discrete system $U' = F(U)$ is integrated with an $s$-stage
Rosenbrock-type method in the transformed form \cite{HairerWanner}
(\code{Integration/Rosenbrock.jl})
\begin{align}
  Y_i &= U^n + \sum_{j<i} a_{ij}\,k_j,\\
  \Bigl(\frac{1}{\gamma\Delta t}I - \mathbf J\Bigr)k_i &= F(Y_i)
    + \frac{1}{\Delta t}\sum_{j<i} c_{ij}\,k_j, \qquad i = 1,\dots,s,\\
  U^{n+1} &= U^n + \sum_{i} m_i\,k_i,
\end{align}
with $A = \alpha\Gamma^{-1}$, $C = -\Gamma^{-1}$ (off-diagonal part) and
$m = \Gamma^{-T}b$. The Jacobian approximation $\mathbf J$ is computed once
per step from the state $U^n$ (W-method) and contains only the vertical
coupling terms (HEVI), so the method is explicit in the horizontal and
linearly implicit in the vertical. The table \code{SSP-Knoth} has
$s=3$, $\gamma=1$,
\begin{equation}
  \alpha = \begin{pmatrix} 0&0&0\\ 1&0&0\\ \tfrac14&\tfrac14&0\end{pmatrix},\quad
  b = \begin{pmatrix}\tfrac16\\ \tfrac16\\ \tfrac23\end{pmatrix},\quad
  \Gamma = \begin{pmatrix} 1&0&0\\ 0&1&0\\ -\tfrac34&-\tfrac34&1\end{pmatrix};
\end{equation}
its explicit part ($\alpha$, $b$) is the three-stage SSP Runge--Kutta method
of Shu and Osher \cite{ShuOsher1988}. Further tables (\code{SSPKnoth12L}, \code{ROS2W},
\code{ROSRK3}, \code{ARS343}, \dots) and explicit Runge--Kutta methods are
available.

\subsection{Vertical Jacobian and Schur complement}

Per column and horizontal node the Jacobian couples $w$ with $\rho$ and
$\rho\theta$ (\code{JacAcousticKernel!}). With the layer thicknesses
$\Delta z_k$, the interface values
$\rho_F = (\rho_k\Delta z_k + \rho_{k+1}\Delta z_{k+1})/(\Delta z_k+\Delta z_{k+1})$
(and analogously $(\rho\theta)_F$) and $\Delta z_{k+1/2} =
\tfrac12(\Delta z_k + \Delta z_{k+1})$, the non-zero blocks are
\begin{align}
  \frac{\partial F_{\rho,k}}{\partial w_{k\pm1/2}} &= \mp\frac{\rho_{F,k\pm1/2}}{\Delta z_k},
  &
  \frac{\partial F_{\rho\theta,k}}{\partial w_{k\pm1/2}} &=
    \mp\frac{(\rho\theta)_{F,k\pm1/2}}{\Delta z_k},\\
  \frac{\partial F_{w,k+1/2}}{\partial (\rho\theta)_{k+1}} &=
    -\frac{\partial_{\rho\theta}p_{k+1}}{\rho_F\,\Delta z_{k+1/2}},
  &
  \frac{\partial F_{w,k+1/2}}{\partial (\rho\theta)_{k}} &=
    \frac{\partial_{\rho\theta}p_{k}}{\rho_F\,\Delta z_{k+1/2}},\\
  \frac{\partial F_{w,k+1/2}}{\partial \rho_{k+1}} &=
    -\frac{g\,\Delta z_{k+1}}{\rho_F(\Delta z_k+\Delta z_{k+1})},
  &
  \frac{\partial F_{w,k+1/2}}{\partial \rho_{k}} &=
    -\frac{g\,\Delta z_{k}}{\rho_F(\Delta z_k+\Delta z_{k+1})}
\end{align}
(plus $\partial_\rho p$ terms, which vanish for the dry $\rho\theta$
formulation). The linear system
\begin{equation}
  \begin{pmatrix}
    \tfrac1{\gamma\Delta t}I & 0 & -J_{\rho w}\\
    0 & \tfrac1{\gamma\Delta t}I & -J_{\rho\theta w}\\
    -J_{w\rho} & -J_{w\rho\theta} & \tfrac1{\gamma\Delta t}I
  \end{pmatrix}
  \begin{pmatrix} k_\rho\\ k_{\rho\theta}\\ k_w\end{pmatrix}
  =
  \begin{pmatrix} r_\rho\\ r_{\rho\theta}\\ r_w\end{pmatrix}
\end{equation}
is reduced to the Schur complement for $w$,
\begin{equation}
  \Bigl(\frac{1}{(\gamma\Delta t)^2}I - J_{w\rho}J_{\rho w}
    - J_{w\rho\theta}J_{\rho\theta w}\Bigr) k_w
  = \frac{1}{\gamma\Delta t} r_w + J_{w\rho} r_\rho + J_{w\rho\theta} r_{\rho\theta},
\end{equation}
which is tridiagonal ($J_{\rho w}$, $J_{\rho\theta w}$ are lower and
$J_{w\rho}$, $J_{w\rho\theta}$ upper bidiagonal). It is solved per column
with the Thomas algorithm (\code{SchurSolveTriKernel!}), followed by the back
substitution $k_\rho = \gamma\Delta t\,(r_\rho + J_{\rho w}k_w)$,
$k_{\rho\theta} = \gamma\Delta t\,(r_{\rho\theta} + J_{\rho\theta w}k_w)$
and $k_{u,v,q} = \gamma\Delta t\,r_{u,v,q}$. The tridiagonal matrix is
assembled once per time step. Optionally the vertical diffusion of the scalars
is also treated implicitly by an extra tridiagonal solve per tracer
(\code{SchurSolveTriDiffKernel!}).

%==============================================================================
\section{Comparison with HOMME and ClimaCore}
%==============================================================================

All three dynamical cores use the same horizontal discretization: continuous
Galerkin spectral elements on the cubed sphere with GLL collocation, lumped
mass, DSS and a vector-invariant momentum equation; they differ in the vertical
coordinate, the vertical discretization, the thermodynamic variable and the
time integration. Table~\ref{tab:comp} summarizes the differences. The
entries for HOMME and ClimaCore are based on the cited literature and
describe the standard configurations.

\begin{table}[h]
\centering\small
\begin{tabularx}{\textwidth}{>{\raggedright\arraybackslash}p{2.6cm}LLL}
\toprule
 & CGDycore.jl (CGSEM) & HOMME / HOMME-NH (E3SM) & ClimaCore / ClimaAtmos (CliMA)\\
\midrule
Horizontal & CG-SEM, GLL, $p=3$--$4$, DSS & CG-SEM, GLL, $n_p=4$ ($p=3$), DSS
  & CG-SEM, GLL, typically $p=3$, DSS\\
Vertical coordinate & terrain-following height & floating Lagrangian
  mass coordinate, remap to hybrid $\eta$ & terrain-following height\\
Vertical discretization & lowest-order staggered, $w$ on interfaces,
  metric from two-point layers
  & mimetic finite differences: $\Delta\pi$, $\Theta$ at mid levels, $w$,
  $\phi$ at interfaces
  & staggered finite differences (centers/faces)\\
Prognostic variables & $\rho$, $u,v$, $w$, $\rho\theta$, $\rho q$
  & $\bm u_h$, $w$, $\phi$, $\Delta\pi$, $\Theta=\theta_v\Delta\pi$, $q\,\Delta\pi$
  & $\rho$, covariant $\bm u_h$, covariant $u_3$, $\rho e_{tot}$, $\rho q$\\
Thermodynamics & potential temperature (energy variants)
  & virtual potential temperature & total energy\\
Vertical transport & limited third-order upwind (Koren type)
  & implicit through the Lagrangian coordinate and conservative remap
  & first/third-order upwind, optional flux-corrected transport\\
Dissipation & biharmonic, separate $\nu$ for curl/div/scalars
  & biharmonic (Laplacian or tensor), subcycled, sponge layer
  & biharmonic hyperdiffusion\\
Tracer limiter & optional optimization-based limiter
  & optimization-based limiter, or semi-Lagrangian transport
  & limiters / flux correction\\
Time integration & Rosenbrock W-methods (HEVI)
  & IMEX Runge--Kutta (HEVI)
  & IMEX additive Runge--Kutta (e.g.\ ARS343)\\
Implicit part & linear vertical acoustics and buoyancy, Schur complement
  for $w$, frozen Jacobian
  & vertical acoustics in $(w,\phi)$, Newton iteration
  & vertical acoustics (and optionally diffusion), Newton with approximate
  Jacobian\\
\bottomrule
\end{tabularx}
\caption{Comparison of the three spectral element dynamical cores.}
\label{tab:comp}
\end{table}

\paragraph{HOMME.}
HOMME (the High-Order Methods Modeling Environment, dynamical core of CAM-SE
and E3SM) introduced the CG-SEM on the cubed sphere for atmospheric models
\cite{TaylorFournier2010,Dennis2012}. The nonhydrostatic version
\cite{Taylor2020} uses a floating Lagrangian vertical coordinate in which the
vertical mass flux vanishes; the layers are remapped conservatively to a
reference hybrid coordinate after a number of steps. The vertical
discretization is mimetic and the semi-discrete scheme conserves the total
energy. Time integration is by IMEX Runge--Kutta methods \cite{Steyer2019}
in which only the vertical acoustic terms are implicit; the resulting
nonlinear system for $w$ and $\phi$ is solved by Newton's method with
tridiagonal linear systems. Hyperviscosity and the limiter
\cite{Guba2014a,Guba2014b} are applied in separate steps.

\paragraph{ClimaCore.}
ClimaCore.jl provides the operators of the CliMA atmosphere model
\cite{Yatunin2025}. It combines the horizontal SEM with a staggered
finite-difference discretization in a terrain-following height coordinate,
like CGDycore. The momentum is stored in covariant components (horizontal
on the cell centers, vertical on the faces) and the thermodynamic variable
is the total energy. The IMEX additive Runge--Kutta methods treat the vertical
acoustic terms implicitly with a Newton iteration and an approximate Jacobian.

\paragraph{Specific features of CGDycore.}
Compared to the two other cores, the CGSEM scheme
\begin{itemize}
\item stores the velocity in local Cartesian (east, north, radial)
  components instead of covariant components, and evaluates all metric terms
  in the conservative (adjugate) form \eqref{eq:metricops} at the two vertical
  nodes of each layer;
\item assembles the vertical velocity with density-weighted half-cell masses,
  which gives a $\bar\rho\,\partial_t w$ form of the vertical momentum
  equation and an exact discrete hydrostatic balance;
\item uses a split (skew-symmetric) form for the horizontal advection of
  $\theta$ and a limited third-order upwind scheme in the vertical;
\item uses linearly implicit Rosenbrock W-methods instead of IMEX Runge--Kutta
  methods: no Newton iteration is needed, the Jacobian is frozen over the
  step and the linear algebra reduces to one tridiagonal solve per column
  and stage;
\item computes the vertical pressure derivative in the terrain-following
  correction of the horizontal pressure gradient from the hydrostatic
  relation.
\end{itemize}

%==============================================================================
\section{Remarks}
%==============================================================================

\begin{itemize}
\item The approximation $\partial_z p \approx -\rho g$ in the terrain term of
  the horizontal pressure gradient (as in MPAS \cite{KlempSkamarock2021}) is
  exact for hydrostatic states. It neglects the vertical derivative of the
  nonhydrostatic pressure perturbation.
\item The kinetic energy and the interface velocity are defined as in
  Section~\ref{sec:kin}. The vertical kinetic energy gradient in the $w$
  equation is assembled half cell by half cell rather than from the layer
  mean $K_k$. This differs from the energy-conserving form by a term
  proportional to $\rho_k - \rho_{k+1}$.
\item The Jacobian contains neither vertical advection nor the
  $u,v$-dependence of the lower boundary condition. Its accuracy influences
  only the stability and the order of the W-method, not the consistency of
  the scheme.
\item The interface treatment of the vertical derivatives in the rotational
  term and in the kinetic energy gradient distributes the jump equally to the
  two adjacent half cells (central treatment). The vertical upwinding only
  acts on the scalar transport.
\end{itemize}

\begin{thebibliography}{99}
\bibitem{TaylorFournier2010}
M.~A. Taylor and A.~Fournier,
A compatible and conservative spectral element method on unstructured grids,
\emph{J. Comput. Phys.} 229 (2010) 5879--5895.

\bibitem{Dennis2012}
J.~M. Dennis, J.~Edwards, K.~J. Evans, O.~Guba, P.~H. Lauritzen,
A.~A. Mirin, A.~St-Cyr, M.~A. Taylor, P.~H. Worley,
CAM-SE: A scalable spectral element dynamical core for the Community
Atmosphere Model,
\emph{Int. J. High Perform. Comput. Appl.} 26 (2012) 74--89.

\bibitem{Taylor2020}
M.~A. Taylor, O.~Guba, A.~Steyer, P.~A. Ullrich, D.~M. Hall, C.~Eldred,
An energy consistent discretization of the nonhydrostatic equations in
primitive variables,
\emph{J. Adv. Model. Earth Syst.} 12 (2020) e2019MS001783.

\bibitem{Steyer2019}
A.~Steyer, C.~J. Vogl, M.~Taylor, O.~Guba,
Efficient IMEX Runge--Kutta methods for nonhydrostatic dynamics,
arXiv:1906.07219 (2019).

\bibitem{Guba2014a}
O.~Guba, M.~A. Taylor, P.~A. Ullrich, J.~R. Overfelt, M.~N. Levy,
The spectral element method (SEM) on variable-resolution grids: evaluating
grid sensitivity and resolution-aware numerical viscosity,
\emph{Geosci. Model Dev.} 7 (2014) 2803--2816.

\bibitem{Guba2014b}
O.~Guba, M.~Taylor, A.~St-Cyr,
Optimization-based limiters for the spectral element method,
\emph{J. Comput. Phys.} 267 (2014) 176--195.

\bibitem{Yatunin2025}
D.~Yatunin, S.~Byrne, C.~Kawczynski, S.~Kandala, G.~Bozzola, A.~Sridhar,
Z.~Shen, A.~Jaruga, J.~Solomon, C.~Kao, et al.,
The CliMA atmosphere dynamical core: concepts, numerics, and scaling,
preprint (2025).

\bibitem{KlempSkamarock2021}
J.~B. Klemp and W.~C. Skamarock,
Adapting the MPAS dynamical core for applications extending into the
thermosphere,
\emph{J. Adv. Model. Earth Syst.} 13 (2021) e2021MS002499.

\bibitem{Koren1993}
B.~Koren,
A robust upwind discretization method for advection, diffusion and source
terms,
in: \emph{Numerical Methods for Advection--Diffusion Problems},
Vieweg, Braunschweig, 1993, pp.~117--138.

\bibitem{HairerWanner}
E.~Hairer and G.~Wanner,
\emph{Solving Ordinary Differential Equations II: Stiff and
Differential-Algebraic Problems}, 2nd ed., Springer, 1996.

\bibitem{ShuOsher1988}
C.-W. Shu and S.~Osher,
Efficient implementation of essentially non-oscillatory shock-capturing
schemes,
\emph{J. Comput. Phys.} 77 (1988) 439--471.
\end{thebibliography}

\end{document}
